
\section{Introduction} 
Vision-Language Models (VLMs) have experienced rapid advancements, demonstrating impressive capabilities in understanding and generating visual content \citep{team2023gemini, achiam2023gpt, li2023blip, dubey2024llama}. These models offer a wide range of functionalities, including image caption generation, visual question answering (VQA), Visual Dialogue, Image Editing, image generation, etc. Examples of such advancements include: (i) Conversation models like Gemini \cite{team2023gemini} and GPT-4o \citep{achiam2023gpt} exhibit strong long-context understanding across image and text modalities, allowing them to analyze complex visual scenes and answer nuanced questions that require reasoning over extended visual and textual information. (ii) Image generation models like Stable Diffusion \citep{rombach2022high}, Imagen \citep{saharia2022photorealistic}, MidJourney, DALL-E\citep{ramesh2021zero}, etc have democratized the creation of highly realistic and diverse visual content from textual prompts. Their increasing accessibility and ease of use empower a wide range of users to generate imagery with unprecedented fidelity and creative control.

The increasing prevalence and capabilities of VLMs increases the criticality of robust safety mechanisms for VLMs across both input and output. For VLMs that accept image inputs, whether synthetic or natural images, it is crucial to build safeguards that prevent harmful content from surfacing. For image generation models, it is crucial to verify compliance with safety policies, preventing the generation of harmful or inappropriate content. This dual challenge underscores the urgent need for highly effective image safety classifiers capable of handling both natural and synthetic images.


The field of image classification has undergone significant transformation with the advent of transformer-based architectures. For example, the Vision Transformer (ViT) \citep{dosovitskiy2020image} processes an image by dividing it into non-overlapping patches, flattening them into sequences, and feeding them into a standard transformer encoder. The Swin Transformer \citep{liu2021swin} introduces a hierarchical structure and a shifted window mechanism to enhance efficiency and scalability while preserving locality. Extending beyond traditional image classification, VLMs such as Gemini, GPT-4o, and Llava have emerged as powerful tools for more comprehensive image understanding tasks, leveraging their ability to process and reason across both visual and textual modalities. However, their direct applicability to specialized vertical domains like image safety classification faces several limitations such as being non-open-sourced, too large and expensive for vertical applications like safety, and not being specifically designed for safety tasks. To bridge this performance gap, recent research has focused on fine-tuning VLMs specifically for image safety classification. Examples include LlavaGuard \citep{helff2024llavaguard} and PerspectiveVision \citep{qu2024unsafebench}, achieving notable improvements.


Despite these advancements, several key limitations remain: (i) Synthetic Data Generation Bottleneck: Existing models often lack automated and targeted training data generation methods. Ideally, a system should be able to produce synthetic images that specifically probe safety boundaries relevant to a particular policy, topic or application. Current approaches often rely on general-purpose datasets that may not adequately cover the diverse and adversarial scenarios necessary for robust safety classification. (ii) Lack of Threshold Customization: some of the existing safety classifiers only provide binary classifications (safe/unsafe) without offering customizable thresholds. Different applications have varying risk tolerances, and the ability to adjust the classification threshold is crucial for balancing precision and recall. 

To address these limitations, we propose ShieldGemma 2 (SG2), a robust image safety classifier fine-tuned on top of the Gemma 3 4B model \cite{gemma_2025}. SG2 offers the following key advantages: 
\begin{itemize}[leftmargin=10pt]
    \item Policy-Aware Classification: SG2 accepts both a user-defined safety policy and an image as input, providing classifications for both natural and synthetic images, tailored to the specific policy guidelines.
    \item Novel Adversarial Synthetic Data Generation: We introduce a novel method for generating synthetic images that are both diverse and adversarial, specifically designed to challenge the classifier based on the needs of the target application. This method ensures more thorough testing and training across a wider range of potential safety violations.
    \item State-of-the-Art Performance (SoTA) with Flexible Thresholding: Internal and external evaluations demonstrate that SG2 achieves SoTA performance on our policies, outperforming prominent models such as LLavaGuard 7B, GPT-4o mini,  and Gemma 3. SG2 outputs a continuous confidence score for each prediction, empowering downstream users to dynamically adjust the classification threshold according to their specific use cases and risk management strategies.
\end{itemize}
